\CQUsection{实验与结果分析}

\CQUsubsection{实验设计概述}
本文实验围绕动脉瘤分割与破裂风险预测两个核心任务展开，分别从体素级病灶识别能力与病例级风险判别能力两个层面对所提出方法进行系统评价。分割实验以 CTA 三维体数据为输入，考察模型输出概率图与人工标注掩膜之间的一致程度；风险预测实验以临床信息、影像表型信息及分割衍生特征为输入，考察模型对动脉瘤破裂状态的区分能力。尽管两项任务在输入形式、输出目标与评价粒度上存在差异，但其共同服务于颅内动脉瘤智能分析这一总体目标，即在影像结构识别基础上进一步实现风险评估。

在实验组织方面，本文严格遵循训练集、验证集和测试集相互独立的基本原则。训练集用于模型参数的学习，验证集用于超参数调节、阈值确定和模型筛选，而测试集仅用于最终性能评估，从而最大限度地减少信息泄漏带来的偏差。
对于分割任务，本文重点考察模型在训练后期的收敛稳定性、概率判别能力和病灶检出能力。对于风险预测任务，则主要对比深度交叉注意力模型与传统机器学习模型在多项评价指标上的表现，并进一步讨论了不同模型的判别特点与适用场景。

\CQUsubsection{实验环境与基本设置}
实验在通用 Linux 计算环境下完成，深度学习模型采用主流张量计算框架进行训练与推理，医学影像读写、空间重采样、数值计算与统计分析均依托成熟的科学计算工具实现。为提高实验过程的可重复性，所有实验均在统一数据划分、统一评价口径和一致训练验证流程下完成记录。本文在结果分析中重点讨论模型表现及其方法学含义，而不涉及具体工程实现细节。

在分割任务中，输入为三维 CTA 体数据。训练阶段采用局部补丁学习策略，以缓解显存限制并提高小目标病灶在训练过程中的可见性；验证阶段则采用滑窗方式进行整例推理，以保证输出结果具有空间完整性。风险预测任务中，输入为病例级表格特征，其中数值变量经标准化处理，类别变量经离散嵌入表示，模型输出为病例级破裂概率。考虑到风险预测任务中正负样本分布通常存在不均衡现象，本文在训练阶段采用类别加权策略，并在验证集上搜索最优判别阈值，以提高模型对少数类样本的识别能力。

为保证比较的公平性，所有对比模型均遵循相同的数据划分原则，测试集不参与训练、调参和阈值选择。分割实验中，模型选择主要依据验证集的重叠指标和阳性检出能力；风险预测实验中，阈值选择仅基于验证集输出概率完成，随后固定该阈值在测试集上报告性能。该设计能够减少因测试集反馈而产生的乐观偏差，使实验结果更能反映模型在未见样本上的泛化表现。

从统计学习角度看，实验设计的关键在于使模型选择过程与最终性能评估过程相互隔离。设训练集、验证集与测试集分别记为 $\mathcal{D}_{tr}$、$\mathcal{D}_{val}$ 和 $\mathcal{D}_{te}$，模型参数由训练集确定：
\begin{equation}
\hat{\theta}=\arg\min_{\theta}\frac{1}{|\mathcal{D}_{tr}|}\sum_{(x_i,y_i)\in\mathcal{D}_{tr}}\mathcal{L}(x_i,y_i;\theta),
\end{equation}
超参数与阈值由验证集确定：
\begin{equation}
\hat{\gamma}=\arg\max_{\gamma\in\Gamma}S_{val}(\hat{\theta},\gamma),
\end{equation}
最终性能则仅在测试集上计算：
\begin{equation}
S_{test}=S(\mathcal{D}_{te};\hat{\theta},\hat{\gamma}).
\end{equation}
其中 $\gamma$ 可表示分类阈值、模型结构选择或其他验证集选择变量。上述形式化描述强调，测试集不应参与任何模型选择过程，否则容易造成性能估计偏乐观。

此外，医学数据通常存在病例间异质性、标注不确定性和类别分布不均衡等特点，因此单一指标难以完整刻画模型行为。本文在分割任务中同时关注区域重叠、阳性检出和概率排序，在风险预测任务中同时关注总体判别、阳性识别和阈值敏感性，以减少由单一评价标准带来的解释偏差。

\CQUsubsection{评价指标}
\CQUsubsubsection{分割任务评价指标}
分割任务的评价需要同时考虑概率判别能力、区域重叠程度与阳性体素检出能力。设真实前景集合为 $Y$，预测前景集合为 $\hat{Y}$，则 Dice 系数定义为
\begin{equation}
\mathrm{Dice}=\frac{2|Y\cap \hat{Y}|}{|Y|+|\hat{Y}|}.
\end{equation}
Dice 系数越高，说明预测区域与真实区域的重叠程度越高。由于动脉瘤通常属于体积较小的局部结构，少量边界偏差便可能引起 Dice 的明显波动，因此该指标对分割质量具有较强敏感性。

在体素级二分类视角下，真阳性、假阳性、真阴性与假阴性分别记为 $\mathrm{TP}$、$\mathrm{FP}$、$\mathrm{TN}$ 和 $\mathrm{FN}$。召回率定义为
\begin{equation}
\mathrm{Recall}=\frac{\mathrm{TP}}{\mathrm{TP}+\mathrm{FN}},
\end{equation}
用于度量真实前景体素中被正确识别的比例。对于医学筛查场景而言，召回率具有重要意义，因为漏检病灶可能直接影响后续临床判断。精确率定义为
\begin{equation}
\mathrm{Precision}=\frac{\mathrm{TP}}{\mathrm{TP}+\mathrm{FP}},
\end{equation}
用于衡量预测为前景的体素中真实属于前景的比例。进一步地，Precision 与 Recall 的调和平均可写为
\begin{equation}
F_1=\frac{2\times \mathrm{Precision}\times \mathrm{Recall}}{\mathrm{Precision}+\mathrm{Recall}}.
\end{equation}
AUC 则从阈值无关角度描述模型对前景与背景体素的整体排序能力，可用于刻画概率输出的可分性。若以真阳性率 $\mathrm{TPR}(\tau)$ 和假阳性率 $\mathrm{FPR}(\tau)$ 分别表示阈值 $\tau$ 下的敏感度和误报率，则 ROC-AUC 可形式化表示为
\begin{equation}
\mathrm{AUC}=\int_{0}^{1}\mathrm{TPR}\left(\mathrm{FPR}^{-1}(u)\right)du,
\end{equation}
其含义是在所有可能阈值下综合评估模型将前景体素排序在背景体素之前的能力。

对于边界敏感的分割评估，还可引入 Hausdorff 距离等轮廓指标。设预测边界点集与真实边界点集分别为 $A$ 和 $B$，则 Hausdorff 距离定义为
\begin{equation}
H(A,B)=\max\left\{\sup_{a\in A}\inf_{b\in B}\|a-b\|,\sup_{b\in B}\inf_{a\in A}\|b-a\|\right\}.
\end{equation}
该指标反映两个边界点集之间的最大偏离程度，适合补充评价分割轮廓的局部极端误差。考虑到小目标分割中边界少量偏差会造成较大指标波动，本文主要以 Dice、Recall、F1 与 AUC 为核心指标，并结合可视化结果讨论边界误差。

\CQUsubsubsection{风险预测评价指标}
风险预测属于病例级不平衡二分类问题，因此仅使用 Accuracy 难以全面反映模型性能。Accuracy 定义为
\begin{equation}
\mathrm{Accuracy}=\frac{\mathrm{TP}+\mathrm{TN}}{\mathrm{TP}+\mathrm{TN}+\mathrm{FP}+\mathrm{FN}},
\end{equation}
反映整体分类正确率，但在正负样本比例失衡时可能掩盖模型对少数类样本的识别不足。因此，本文同时报告 Precision、Recall、F1 与 AUC 等指标，以更全面地刻画模型在风险识别任务中的综合表现。

对于风险预测模型输出的概率 $\hat{p}_i$，其二分类结果由阈值 $\tau$ 决定：
\begin{equation}
\hat{y}_i=
\begin{cases}
1, & \hat{p}_i\geq \tau,\\
0, & \hat{p}_i<\tau.
\end{cases}
\end{equation}
本文在验证集上搜索较优阈值，并将其固定用于测试集评估。该策略能够依据数据分布与任务目标自适应调整决策边界，从而减少默认阈值对实验结论的影响。例如，当以 F1 值作为阈值选择准则时，可写为
\begin{equation}
\tau^{\ast}=\arg\max_{\tau\in[0,1]}\frac{2\mathrm{Precision}(\tau)\mathrm{Recall}(\tau)}{\mathrm{Precision}(\tau)+\mathrm{Recall}(\tau)}.
\end{equation}

除上述指标外，医学风险预测还需要关注概率输出的可解释性和校准性。理想情况下，当模型对一组病例给出接近 $p$ 的风险概率时，该组病例中实际阳性比例也应接近 $p$。若将样本按预测概率划分为 $K$ 个区间 $B_k$，则期望校准误差可写为
\begin{equation}
\mathrm{ECE}=\sum_{k=1}^{K}\frac{|B_k|}{N}\left|\operatorname{acc}(B_k)-\operatorname{conf}(B_k)\right|,
\end{equation}
其中 $\operatorname{acc}(B_k)$ 表示第 $k$ 个区间内的真实阳性比例，$\operatorname{conf}(B_k)$ 表示平均预测概率。虽然本文重点报告分类判别指标，但校准分析对于后续临床使用具有重要意义，因为风险概率不仅影响二分类决策，也影响临床复核和随访优先级安排。

为便于说明各类评价指标在本文中的功能定位，表~\ref{tab:metric_summary} 对本章主要指标及其分析侧重点进行了归纳。

\begin{table}[htbp]
\centering
\caption{主要评价指标及其分析侧重点}
\label{tab:metric_summary}
\small
\begin{tabular}{lll}
\toprule
任务类型 & 指标 & 主要反映内容 \\
\midrule
分割任务 & Dice & 预测区域与真实区域的重叠程度 \\
分割任务 & Recall & 病灶前景体素的检出能力 \\
分割任务 & Precision & 前景预测结果的可靠性 \\
分割任务 & AUC & 前景与背景体素的整体可分性 \\
风险预测任务 & Accuracy & 整体分类正确率 \\
风险预测任务 & Precision & 高风险预测结果的准确性 \\
风险预测任务 & Recall & 破裂病例的识别能力 \\
风险预测任务 & F1 & Precision 与 Recall 的综合平衡 \\
风险预测任务 & AUC & 破裂与未破裂病例的排序能力 \\
\bottomrule
\end{tabular}
\normalsize
\end{table}

\CQUsubsection{动脉瘤分割实验}
\CQUsubsubsection{实验目的与训练过程}

分割实验的主要目的是验证三维 Attention U-Net 及其改进结构对颅内动脉瘤的定位能力和轮廓恢复能力。训练时，模型通过局部补丁样本学习病灶及其邻域的空间结构；在验证和测试阶段，则对整例 CTA 体数据进行滑窗推理。由于动脉瘤在空间上极为稀疏，实验不仅关注病灶前景的召回率，还着重考虑在复杂血管背景下的误检控制。
从训练过程来看，模型在后期的验证指标逐渐趋于稳定，表明网络已基本完成从低层纹理到高层病灶语义的表征学习。本文重点统计了这一阶段的验证结果，以观察收敛末期的性能表现。在此期间，AUC 和 Recall 均保持在相对稳定的区间，说明模型对前景与背景的概率排序能力较强，且病灶检出性能未出现明显波动。
在优化曲线解读上，如果训练损失持续下降而验证指标不再提升，通常意味着模型可能开始拟合训练集中的噪声；若训练损失与验证指标同步改善，则说明模型仍在学习具有泛化价值的特征。因此，本文并未单纯以训练损失作为模型选择的依据，而是综合考虑验证集上的重叠指标、检出指标和概率排序指标进行判断。对于动脉瘤这类小目标任务，验证集 Recall 的稳定性尤为关键，因为较高的 AUC 并不一定能保证所有小病灶在固定阈值下都被正确检出。
此外，训练过程还反映出小目标分割中“概率排序”与“空间成形”之间的区别。前者侧重模型能否给真实病灶体素赋予高于背景的概率，后者则关注阈值化后是否能形成连续、解剖合理的三维区域。若模型 AUC 较高但 Dice 较低，往往说明其概率排序能力已经较好，但阈值化后的分割结果可能存在边界收缩、离散假阳性或局部漏检等问题。因此，在训练后期，除了关注损失函数和 AUC 外，还需要结合固定阈值下的 Dice、Recall 以及可视化结果，综合判断模型是否具备稳定的空间分割能力。


\CQUsubsubsection{分割结果与阈值分析}
以较优验证轮次为代表，在阈值 0.5 下模型取得 AUC 为 0.97、Recall 为 0.56 的结果，表明模型在当前验证数据上已具备较好的体素级区分能力。以最终检查点作为报告对象，在相同阈值下 AUC 为 0.96、Recall 为 0.55，与较优轮次表现接近，说明训练后期模型并未出现明显退化，整体泛化性能并未显著降低。

需要指出的是，较高的 AUC 并不意味着分割结果已经在空间层面达到完全理想的状态。AUC 主要衡量不同阈值下模型对正负体素的整体排序质量，而最终二值掩膜的空间一致性还受到阈值选择、病灶体积、边界清晰度以及重建方式等因素影响。对于小目标病灶而言，即便模型在概率排序层面表现良好，不合适的阈值仍可能导致边界收缩、局部断裂或离散误检区域。因此，对分割模型的评价应结合 Dice、Recall 与典型病例可视化结果综合分析，而不宜依赖单一指标作出判断。

阈值扫描结果表明，随着分割阈值 $\tau_s$ 的提高，模型预测趋于保守，误检体素数量可能减少，但 Recall 往往同步下降；反之，当阈值降低时，模型更倾向于保留低置信度疑似前景，Recall 可能上升，但误检风险也会增加。其决策形式可写为
\begin{equation}
\hat{M}_{i,j,k}(\tau_s)=\mathbb{I}\left(P_{i,j,k}\geq \tau_s\right).
\end{equation}
当 $\tau_s$ 增大时，满足条件的前景集合单调收缩，因此漏检概率通常增加。对于医学辅助分析场景而言，维持较高的病灶检出能力具有重要意义，但同时也需控制误检区域规模，以保证后续结构定量分析与临床展示的可解释性。

\CQUsubsubsection{可视化分析与误差来源}
从典型病例的可视化结果看，模型在多数病灶区域能够较完整地覆盖动脉瘤主体，并在局部边界上形成相对连续的响应。对于形态较规则、对比度较高的病灶，预测掩膜与人工标注具有较高一致性；而当病灶体积较小、边缘对比度不足、与载瘤血管过渡模糊或邻近高密度结构干扰明显时，模型可能出现局部漏检、边界偏移或少量误检。

这类误差具有一定普遍性。一方面，小病灶在多次下采样过程中容易发生空间信息损失，即使通过跳跃连接进行细节补偿，仍可能难以完全恢复极细小结构边界；另一方面，CTA 中血管增强区域与动脉瘤前景在灰度分布上具有相似性，模型需要借助空间上下文与拓扑关系而非单纯强度差异进行判别。由此可见，后续若需进一步提升分割性能，可从更具针对性的前景采样、边界敏感损失设计、多尺度监督及稳健后处理策略等方面展开优化。

从误差分解角度看，分割误差可大致分为阈值误差、边界误差和结构误差三类。阈值误差主要源于概率输出与二值决策之间的不匹配；边界误差主要表现为预测轮廓相对于真实轮廓的局部偏移；结构误差则表现为病灶区域断裂、漏检或与邻近血管粘连。设预测概率为 $P$，真实标签为 $Y$，阈值化结果为 $\hat{M}(\tau_s)$，则不同阈值下的分割质量可表示为
\begin{equation}
S(\tau_s)=\mathcal{M}\left(\hat{M}(\tau_s),Y\right),
\end{equation}
其中 $\mathcal{M}(\cdot)$ 表示任一分割评价指标。若 $S(\tau_s)$ 在较宽阈值区间内保持稳定，说明模型概率图具有较好的校准性与空间鲁棒性；若指标随阈值剧烈波动，则提示模型输出中存在较多低置信度或边界不确定区域。

对于动脉瘤分割而言，误差分析还应关注空间位置和解剖背景。病灶与载瘤血管交界处通常是最容易发生边界偏移的区域，原因在于二者在增强影像中具有相近灰度分布，且真实标注本身也可能受到观察者经验影响。邻近骨性结构、高密度伪影或血管分叉区域则可能诱发假阳性响应。上述现象说明，模型性能改进不能仅依赖更复杂的网络结构，还需要从数据表示、监督信号和评价方式上共同约束。

\CQUsubsubsection{分割模型消融分析}
为进一步分析模型性能提升的来源，本文从结构设计角度对分割模型进行了消融讨论。可将基础三维 Attention U-Net、引入 ASPP 的增强模型以及进一步加入三维坐标注意力的模型视为逐步增强的结构序列。基础模型主要依赖编码器—解码器结构和注意力门控完成病灶恢复；ASPP 模块通过多尺度空洞卷积扩大感受野，有助于增强模型对不同尺度病灶及其邻域上下文的刻画能力；三维坐标注意力则进一步引入方向感知的位置建模，对边界模糊、形态不规则或空间走向复杂的病灶具有潜在优势。

在现有实验结果中，增强模型相较于基础模型表现出更优的综合分割性能。在相同训练轮次条件下，ASPP-CA-UNet 的 Recall 达到 0.67，AUC 达到 0.96；与基础模型相比，其 Dice 系数提升约 20\%，F1 值提升约 40\%。此外，两类模型在绝大多数背景体素上均保持较高一致性，说明增强模型性能提升并非来源于对背景的过度拟合，而更可能归因于对病灶区域表征能力的增强。为便于比较，表~\ref{tab:seg_ablation} 对分割消融结果作了汇总。

\begin{table}[htbp]
\centering
\caption{分割模型消融结果比较}
\label{tab:seg_ablation}
\small
\begin{tabular}{lcccc}
\toprule
模型 & Recall & AUC & Dice（相对变化） & F1（相对变化） \\
\midrule
基础三维 Attention U-Net & 0.56 & 0.97 & 基准水平 & 基准水平 \\
Attention U-Net + ASPP + CA & 0.67 & 0.96 & +20\% & +40\% \\
\bottomrule
\end{tabular}
\normalsize
\end{table}

需要说明的是，由于动脉瘤分割属于高度不平衡的体素级任务，背景体素占比远高于病灶体素，因此“整体像素一致性”往往会受到大量真阴性体素的主导，其解释价值相对有限。相比之下，Dice、Recall 与 F1 等更能体现病灶区域识别能力的指标，能够更加真实地反映模型在本任务中的性能差异。基于此，本文在消融分析中更强调增强结构对病灶区域识别质量的改善，而非单纯关注整体一致性指标。

ASPP 与坐标注意力的作用机制和增益来源有所不同。ASPP 通过不同空洞率的并行分支来扩大感受野，主要提升模型对不同尺度病灶及其周围血管背景的上下文理解能力；而坐标注意力则通过对空间方向信息的建模来强化位置感知，更有利于处理走向复杂或边界模糊的血管结构。
如果只引入多尺度特征而缺乏位置约束，模型容易加强对背景血管的响应；反之，若仅加入位置注意力而缺少多尺度上下文，对形态变化较大的病灶适应性则可能不足。因此，这两类模块在理论上具有较好的互补性，联合使用能更有效地平衡小目标分割中的背景抑制与结构保持问题。


若用 $S_0$ 表示基础模型性能，用 $S_A$ 表示加入 ASPP 后的性能，用 $S_{A+C}$ 表示同时加入 ASPP 与坐标注意力后的性能，则各模块的边际贡献可形式化表示为
\begin{align}
\Delta_{A} &= S_A-S_0,\\
\Delta_{C\mid A} &= S_{A+C}-S_A.
\end{align}
其中 $\Delta_A$ 反映多尺度上下文增强的直接贡献，$\Delta_{C\mid A}$ 反映在已有多尺度表示基础上方向注意力带来的增量效果。若二者均为正，则说明两个模块对目标指标具有累积增益；若单独增益较弱但联合增益明显，则可能存在互补作用。本文实验结果显示，增强结构在 Dice 与 F1 上的相对提升更为明显，表明其主要改善的是前景区域的重叠质量与阳性预测可靠性，而不仅是概率排序能力。

\CQUsubsection{破裂风险预测实验}
\CQUsubsubsection{实验目的与模型设置}
破裂风险预测实验旨在评价多源表格特征对病例级破裂状态的判别能力，并比较交叉注意力模型与传统机器学习模型在中小规模医学数据条件下的表现。所有模型均基于相同训练集、验证集与测试集划分进行评估，并在验证集上完成阈值选择，以保证不同模型之间具有统一的比较基础。

本文纳入比较的模型包括交叉注意力深度表格模型、随机森林、逻辑回归、支持向量机、$k$ 近邻以及融合模型。逻辑回归可视为线性可解释基线，随机森林代表非线性集成树模型，支持向量机与 $k$ 近邻分别代表基于核映射与局部相似性判别的传统分类方法，而交叉注意力模型则代表通过字段交互学习高阶依赖关系的深度方法。融合模型进一步综合树模型与注意力模型输出，用于检验不同模型范式之间的互补性。

\CQUsubsubsection{对比实验结果}
风险预测模型测试集结果如表~\ref{tab:risk_result} 所示。表中 Threshold 表示由验证集确定并固定用于测试集评估的分类阈值。

\begin{table}[htbp]
\centering
\caption{风险预测模型测试集对比结果}
\label{tab:risk_result}
\small
\setlength{\tabcolsep}{4pt}
\resizebox{\textwidth}{!}{%
\begin{tabular}{lcccccc}
\toprule
模型 & Accuracy & Precision & Recall & F1 & AUC & Threshold \\
\midrule
Forest-Attention Ensemble & 0.7944 & 0.6931 & 0.7778 & \textbf{0.7330} & \textbf{0.8568} & 0.350 \\
Risk Cross-Attention & 0.7852 & 0.6832 & 0.7667 & 0.7266 & 0.8546 & 0.325 \\
Random Forest & 0.7419 & 0.6182 & 0.7556 & 0.6800 & 0.8284 & 0.400 \\
Logistic Regression & 0.6895 & 0.5512 & 0.7778 & 0.6452 & 0.7342 & 0.450 \\
KNN & 0.5766 & 0.4576 & \textbf{0.9000} & 0.6067 & 0.6966 & 0.200 \\
SVM-RBF & 0.5565 & 0.4451 & \textbf{0.9000} & 0.5956 & 0.6732 & 0.300 \\
\bottomrule
\end{tabular}%
}
\normalsize
\end{table}

从表~\ref{tab:risk_result} 可以看出，Forest-Attention Ensemble 在测试集上取得最高的 F1 值 0.7330 与 AUC 值 0.8568，说明传统树模型与深度注意力模型在当前任务中具有一定互补优势。Risk Cross-Attention 模型取得 F1 值 0.7266 和 AUC 值 0.8546，其性能与最优融合模型非常接近，并明显优于逻辑回归、支持向量机及 $k$ 近邻等传统单模型基线。这表明基于字段级交互的深度建模方式能够较有效地捕获临床变量与影像表型变量之间的非线性关联。

进一步分析可见，随机森林在 Accuracy、F1 和 AUC 等多项指标上均优于多个传统基线，体现出集成树模型在中小规模表格数据上的稳定性与适应性。相比之下的KNN 与 SVM-RBF 模型虽然在 Recall 上达到 0.9000，表现出较强的阳性病例检出倾向，但其 Precision 分别仅为 0.4576 和 0.4451，同时 F1 与 AUC 均处于较低水平。这说明上述模型更倾向于将样本判定为高风险类别，在此虽略微提高了召回能力，但是却显著增加了误报率，从整体上而言并未取得超过随机森林或随机森林-注意力融合模型的结果。逻辑回归具有较强可解释性，但从其 AUC 仅为 0.7342可见线性决策边界难以充分刻画当前多源特征中的复杂非线性关系。

从模型假设角度看，不同方法的差异可归纳为判别函数空间的差异。逻辑回归对应线性形式
\begin{equation}
\log\frac{\hat{p}}{1-\hat{p}}=\beta_0+\sum_{j}\beta_jx_j,
\end{equation}
其优势在于参数含义明确，但难以自动表达高阶非线性交互。随机森林通过多棵决策树的集成形成分段非线性判别函数，可较好处理变量间的条件分裂关系。交叉注意力模型则通过字段 token 之间的注意力权重学习全局交互，其判别函数可抽象为
\begin{equation}
\hat{p}=\sigma\left(F_{\omega}(e_1,e_2,\ldots,e_{m+q})\right),
\end{equation}
其中 $F_{\omega}(\cdot)$ 表示由多层注意力与非线性映射构成的可学习函数。该类模型具有更强表达能力，但对样本规模、正则化和数据噪声更敏感。因此，实验结果中深度模型与集成树模型表现接近，符合中小规模医学表格数据的一般规律。

\CQUsubsubsection{阈值选择与不平衡分类讨论}

风险预测实验中，各模型的最优阈值普遍低于或接近 0.5，这表明在类别不平衡条件下，默认判别阈值未必能够适配医学风险识别任务的实际需求。对于输出概率 $\hat{p}$ 而言，降低阈值会使更多病例被判定为高风险，进而提高 Recall，但也可能降低 Precision；提高阈值则会产生相反效果。因此，阈值选择本质上是在漏诊风险与误报代价之间进行权衡。

从临床辅助分析角度看，破裂风险预测通常更关注高风险病例的检出能力，因此相对其他属性而言， Recall 具有较高参考价值；但若 Precision 过低，则可能导致过多低风险病例被误判为高风险，从而增加后续检查、复核与干预负担。F1 值通过综合考察 Precision 与 Recall，为动脉瘤破裂风险预测的二分类任务提供了更平衡的评价方案。AUC 则从全阈值范围反映模型对阳性与阴性病例的排序能力，可用于模型间总体判别能力的比较。

进一步地，阈值选择可被视为代价敏感决策问题。若假阴性代价为 $C_{FN}$，假阳性代价为 $C_{FP}$，则在给定预测概率 $\hat{p}$ 时，将病例判定为阳性的期望代价为 $C_{FP}(1-\hat{p})$，判定为阴性的期望代价为 $C_{FN}\hat{p}$。当
\begin{equation}
C_{FP}(1-\hat{p}) < C_{FN}\hat{p}
\end{equation}
时，选择阳性判定更为合理。由此可得理论阈值
\begin{equation}
\tau_c=\frac{C_{FP}}{C_{FP}+C_{FN}}.
\end{equation}
当漏诊代价高于误报代价时，$C_{FN}>C_{FP}$，对应阈值 $\tau_c$ 将低于 0.5，这与实验中多个模型的最优阈值低于默认阈值的现象一致。该分析表明，阈值并非纯粹的技术参数，而反映了具体应用场景下对漏诊与误报的权衡。

\CQUsubsubsection{风险模型结果讨论}

综合各项指标来看，基于交叉注意力的深度表格模型在本次任务中表现出了较强的竞争力，但并未全面超越所有传统方法。其主要原因在于：一方面，交叉注意力机制能够显式捕捉异构字段间的复杂依赖关系，在处理多源医学表格数据时具有较好的表达能力；另一方面，在样本规模相对有限的情况下，传统集成学习模型通常表现出更好的稳健性和更低的过拟合风险。

因此，深度模型与传统模型之间并非简单的替代关系，而是在不同数据条件和评价目标下各有优势。在融合模型在本实验中取得了最优的综合性能的这一结果进一步验证了这一判断。由此表明，基于树模型的分裂规则与基于注意力机制的表示学习其实捕捉到了不同类型的判别模式，二者结合能在一定程度上缓解单一模型的偏差，从而提升整体分类效果。这也为后续研究探索更系统的模型融合策略提供了支持，从而进一步提高风险预测的稳定性和泛化能力。

从误差类型来看，假阴性和假阳性在风险预测任务中具有不同的临床意义。假阴性意味着高风险病例被低估，可能延误复查和干预时机；假阳性则意味着低风险病例被高估，可能会增加不必要的随访负担和临床资源消耗。因此，在该任务中不能单纯追求 Accuracy，而应根据具体临床场景权衡 Recall、Precision 和 F1 等指标的相对重要性。本实验中，部分模型通过降低阈值获得了较高的 Recall，但 Precision 明显下降，说明其对阳性类别的概率输出较为敏感，却缺乏足够的可靠性。相比之下，交叉注意力模型和融合模型在 Precision 与 Recall 之间取得了更为均衡的表现，综合评价指标也更高。

同时，当前结果也提示应谨慎理解深度表格模型的优势。在样本数量有限、变量噪声较高或阳性比例不足时，复杂模型可能学习到与训练分布相关但泛化能力有限的模式。因此，较高的测试集指标需要结合置信区间、重复划分或外部验证进一步确认。本文实验结果表明交叉注意力模型具备较好的风险判别潜力，但其临床稳定性仍需在更大规模数据和更严格验证条件下进一步评估。

\CQUsubsection{结果综合讨论}

综合分割和风险预测实验结果可以看出，本文方法在当前数据集上初步构建了从影像结构识别到病例风险估计的完整分析流程。分割模型在体素级 AUC 和 Recall 上表现稳定，表明其能够较好地区分病灶与复杂血管背景；风险预测模型在 F1 和 AUC 上接近最优融合模型，显示交叉注意力机制在多源表格特征建模中具有较好的适用性。

从方法构成来看，分割任务通过注意力门控、ASPP 和三维坐标注意力分别实现了背景抑制、多尺度上下文聚合与方向位置建模，三者共同提升了模型对复杂血管结构的表达能力。而在风险预测任务中，字段级嵌入与交叉注意力机制将传统的特征拼接转化为可学习的交互表示，使模型能够在病例层面捕捉临床变量、影像表型和形态特征之间的潜在依赖关系。

同时，实验也暴露了一些值得关注的问题。首先，医学影像数据的规模有限，标注成本高，单次数据划分结果容易受样本分布影响；其次，风险预测任务中的类别不平衡问题会显著影响阈值选择和指标解读；此外，深度表格模型在中小样本医学任务中并不总是稳定优于传统集成模型，合理的正则化、特征选择、概率校准以及外部验证仍非常必要。

\CQUsubsection{系统闭环验证}
为验证整体方法流程的完整性，本文将分割推理、特征提取与风险预测组织为端到端演示闭环。该闭环能够从输入影像出发，依次完成病灶概率图生成、二值掩膜输出、结构信息汇聚与破裂风险概率估计，并以可视化形式呈现关键结果。系统闭环验证的意义不在于替代临床诊断，而在于证明所提出的分割方法与风险建模方法可以在统一分析流程中协同运行，从而为后续临床可用性研究与系统化部署提供原型基础。

\CQUsubsection{本章小结}
本章从实验设计、评价指标、分割实验、风险预测对比、消融分析与结果讨论等方面，对本文方法进行了较为系统的实验性评价。结果表明，分割模型在当前验证数据上具有较好的体素级判别能力与病灶检出能力；风险预测模型能够有效利用多源表格特征，并在测试集上取得接近最优融合模型的综合表现。总体而言，本文方法在病灶识别与风险评估两个层面均达到了预期目标，但其稳健性，泛化能力的检测与提升仍有赖于更大规模、多中心的多轮重复实验的进一步验证与优化。